\chapter{Framework GamiDOC}
\label{ch:fgamidoc}
Il GamiDOC (Bassanelli) è il documento di game design richiesto tra i 5
deliverable d'esame. Le sezioni ufficiali sono:

\begin{enumerate}
\item \textbf{Context}: pubblico, bisogni, vincoli didattici (roadmap Dart/Flutter).
\item \textbf{Game behavior}: loop Studio (6--8\,min) $\rightarrow$ Quiz
(3--5 domande) $\rightarrow$ Deck (2--3 carte) $\rightarrow$ Boss (10--15\,min).
\item \textbf{Gamification}: punti, badge, level, collezioni, regole di
sblocco (soglia 80\%, tentativi illimitati ma penalizzati, energia 3+, vite 3).
\item \textbf{Aesthetics}: fantasy dungeon (Gargoyle, StackOverflow, Async,
WidgetTitan), rarità Common 6 / Rare 10 / Epic 16+Perf5 / Legendary 25+Crit25\%.
\item \textbf{Dynamics}: Dual Progression (XP giocatore + potere deck).
\item \textbf{Architettura}: rimando tecnico (si veda \texttt{specifica.pdf}).
\item \textbf{Valutazione}: protocollo utente da eseguire (si veda il
Cap.~\ref{ch:userresearch}).
\item \textbf{KPI/Analytics}: da definire (retention, mastery, tempo/boss).
\end{enumerate}

Il GamiDOC completo e presentabile è \texttt{gamidoc.pdf}; qui ne fissiamo
la struttura teorica, lì stanno i contenuti di progetto.
